\documentclass[12pt, letterpaper]{report}
\usepackage{graphicx}
\usepackage{amsmath}

% Global equation numbering (not by chapter)
\counterwithout{equation}{chapter}

% Do the same for figures and tables if you want
\counterwithout{figure}{chapter}
\counterwithout{table}{chapter}

\def\mathbi#1{\textbf{\em #1}}

\graphicspath{{../src/workspace/outputs_2D/plots/sparse/}}

\title{Inferred solution of Unsteady Stokes Flow using Physics Informed Gaussian Process}
\author{Mingkwan RATTANASIWAMOK}
\date{September 2026}
\begin{document}
\maketitle


\pagebreak  % Preferred - allows LaTeX to adjust spacing

\chapter{Introduction}
    \section{Unsteady Stokes FLow}
Behavior of fluid could be explain via the Navier Stokes Equation, which utilize Cauchy Momentum to do so. For an imcompressible fluid,

\begin{equation}
    \rho \frac{\partial u}{\partial t} + \nabla (u \cdot \nabla) = \nabla p + \eta{\nabla}^{2} u + f_{ext}
\end{equation}

For fluid with low Reynold number, the self-advection term could be neglected 

\begin{equation}
    \rho \frac{\partial u}{\partial t} = -\nabla p + { \eta\nabla}^{2} u + f_{ext}
\end{equation}

Unsteady Stokes flow is blah blah

\section{Gaussian Process Regression}

Gaussian Process or GP are defined by a series of infinite value that have a joint Gaussian distribution.
A GP for function \(f(x)\), with the input locations at \(\mathbf{X} =[\mathbf{x_1},\mathbf{x_2},...,\mathbf{x_N}] \), will provide values for
\(\mathbf{f(X)}=[\mathbf{f(x_1)},\mathbf{f(x_2)},...,\mathbf{f(x_N)}]\). The joint distribution is in the form of a multi-variate Gaussian

\begin{equation}
    p(\mathbf{f(X)} | \mathbf{X},\mathbf{\theta},\mu(\cdot),k(\cdot,\cdot)) = \mathcal{N}\left(
        \mu(\mathbf{X},k\mathbf(X,X;\theta))
    \right)
\end{equation}

\(\mu\) and \(k\) represent the average function and kernel function resepectively

\begin{equation}
    \mu(\mathbf{X})=\langle f(\mathbf{X}) \rangle \equiv \langle f \rangle (\mathbf(x))
\end{equation}

\begin{equation}
    k\mathbf(X,X;\theta)=\langle \delta f(\mathbf{x}) \delta f(\mathbf{x^{'}}) \rangle \equiv \langle \delta f\delta f (\mathbf{x},\mathbf{x^{'}}) \rangle
\end{equation}

where \(\theta\) is the hyperparameter of the kernel function. The kernel function used was a square expotential kernel.

\begin{equation}
    k(x,x^{'})=exp \left( -\frac{(x-x^{'})^2}{2l^2} + \gamma\right)
\end{equation}

to capture the dimensionality of the problem, derived kernels were utilized such as product kernel, to calculate the kernel function.

\begin{equation}
    k^{\times}(\mathbf{x},\mathbf{x^{'}})=k(x,x^{'}) \times k(y,y^{'})
\end{equation}


We can learn the function by inferring the value at each test positions \(X_{\star}\). For a function \(f\),

\begin{equation}
    \begin{bmatrix}
        \mathbf{f^{\star}} \\
        \mathbf{f} \\
    \end{bmatrix}
        \sim
    \mathcal{N}\left(
        \begin{bmatrix}
            \mathbf{\mu^{\star}} \\
            \mathbf{\mu} \\
        \end{bmatrix}
        ,
        \begin{bmatrix}
            \mathbf{K_{\star \star}} & \mathbf{K_{\star}}\\
                & \mathbf{K} \\
        \end{bmatrix}
    \right)
\end{equation}

the distribution of \(\mathbf{f_{\star}}\ | \mathbf{f} 
    \sim
    \mathcal{N}\left(
        \mathbf{\mu} + \mathbf{K_{\star}}\mathbf{K^{-1}}\delta \mathbf{f} ,  \mathbf{K_{\star \star}} - \mathbf{K_{\star}}\mathbf{K^{-1}}\mathbf{K^{T}_{\star}}
    \right)\) 


\pagebreak

\section{Integration of Physics}

    \subsection{Main Idea}

Using backward Euler we can express \(u^t\) as following

\begin{equation}
    u^t = u^{t-1} + \frac{\Delta t}{\rho} \left({-\nabla p +  \eta{\nabla}^{2} u + f_{ext}}\right)
\end{equation}

Now, 

\begin{equation}
    u^t + \frac{\Delta t}{\rho} \left({\nabla p -  \eta{\nabla}^{2} u}\right)  = u^{t-1} + f_{ext}
\end{equation}

The \(LHS\) of the equation, we defined as \(\mathcal{G}\)

\begin{equation}
    \mathcal{G} \equiv u^t + \frac{\Delta t}{\rho} \left({\nabla p -  \eta{\nabla}^{2} u}\right)
\end{equation}
\[\mathcal{G} \equiv u^t + \frac{\Delta t}{\rho} f\]
\[\mathcal{G} \equiv u^{t-1} + f_{ext}\]

where \(f\) is 

\begin{equation}
    f = {\nabla p - \eta{\nabla}^{2} u}
\end{equation}

\subsection{Prior} 

We can learn the function by inferring the value at each test positions \(X_{\star}\). For a function \(f\),

\begin{equation}
    \begin{pmatrix}
        \mathbf{u} \\
        \mathbf{p} \\
        \mathbf{\mathcal{G}} \\
        \mathbf{s} \\
    \end{pmatrix}
    \sim
    \mathcal{N}\left(
        0, \mathbf{K}
    \right)
\end{equation}

\begin{equation}
    K = 
    \begin{bmatrix}
        \mathbf{K_{u u}} &\mathbf{K_{u p}}& \mathbf{K_{u G}}& \mathbf{K_{u s}} \\
        \ & \mathbf{K_{p p}}& \mathbf{K_{p G}}& \mathbf{K_{p s}} \\
        \ & & \mathbf{K_{G G}} & \mathbf{K_{G s}} \\
        \ & & & \mathbf{K_{s s}} \\
    \end{bmatrix}
\end{equation}

    \subsection{Posterior} 

The prior of all the observed points and governing equation can be expressed as

\begin{equation}
\begin{bmatrix}
u_b\\
S^1_L \\
S^1_p \\
\mathcal{G}^{t} \\
s
\end{bmatrix}
\sim\mathcal{N}\left(0,
\begin{bmatrix}
\mathbf{K_{\star \star}} & \mathbf{Q}\\
\mathbf{Q^T} & \mathbf{K} \\
\end{bmatrix}
\right)
\end{equation}
where, \(u_b\) denote the observed boundary value of the velocity.
The correlation matrixes can be writtten as

\begin{equation}
\begin{bmatrix}
\mathbf{K_{u u}} &\mathbf{K_{u p}}& \mathbf{K_{u G}}& \mathbf{K_{u s}} \\
\ & \mathbf{K_{p p}}& \mathbf{K_{p G}}& \mathbf{K_{p s}} \\
\ & & \mathbf{K_{G G}} & \mathbf{K_{G s}} \\
\ & & & \mathbf{K_{s s}} \\
\end{bmatrix}
\end{equation}

The full posterior can be expressed as following, where \(\mathcal{G}^{t-1}\) is defined as the \emph{RHS} of the unsteady stokes equation.

\[
\begin{bmatrix}
u^{t}\\
p^{t} \\
\end{bmatrix}
|
\begin{bmatrix}
u_b\\
S^1_L \\
S^1_p \\
\mathcal{G}^{t} \\
s
\end{bmatrix}
\sim\mathcal{N}\left(\mathbf{q^{T}}\right)
\]


\chapter{Methodology}

\chapter{Results}

1. dt = 0.05 and 5/10 loops work

2. dt = 0.01 and 25 loops work (Testing the effect of big and small time interval) {Should be fine}

2. dt = 0.01 and 25 and 200, 270, 340 datapoints work (Testing the effect training artificial data sizre) {Should be fine}

Note' right now the optimizer is optimizing the same function on the same value. At the start of each loop, the data should reset so a new hyperparam could be train on new data.

\end{document}